% Generated draft: requires mathematical and source review.
\documentclass[11pt]{article}
\usepackage[margin=25mm]{geometry}
\usepackage{fontspec}
\setmainfont{Latin Modern Roman}
\usepackage{amsmath,amssymb,mathtools,unicode-math}
\setmathfont{Latin Modern Math}
\usepackage{enumitem,needspace,xurl,hyperref}
\setlength{\parindent}{0pt}
\setlength{\parskip}{5pt}
\setlength{\emergencystretch}{3em}
\newcommand{\N}{\mathbb{N}}
\newcommand{\Z}{\mathbb{Z}}
\newcommand{\Q}{\mathbb{Q}}
\newcommand{\R}{\mathbb{R}}
\newcommand{\C}{\mathbb{C}}
\newcommand{\F}{\mathbb{F}}
\newcommand{\A}{\mathbb{A}}
\newcommand{\PP}{\mathbb{P}}
\newcommand{\E}{\mathbb{E}}
\newcommand{\eps}{\varepsilon}
\newcommand{\dd}{\,\mathrm{d}}
\begin{document}
\section*{UnsolvedMath ID: 700887}
\textbf{Source problem code:} catalogue entry
\subsection{Definitions and mathematical statement}
Let \(F\) be a countably infinite field, let \(G\leq\operatorname{Aut}(F)\) be a finite group, and let \(S\subseteq F\) be a subfield. The integral group ring is
\[
\Z G=\left\{\sum_{g\in G}n_g g:n_g\in\Z\right\};
\]
an element is nonzero precisely when its coefficients are not all zero. The source imposes the condition
\[
\forall x\in F\quad\exists f\in\Z G\setminus\{0\}\quad x^f\in S,
\]
and asks whether this implies that \(F=S\) or that \(F/S\) is purely inseparable. The element \(f\) may depend on \(x\), and no invariance of \(S\) under \(G\) is assumed.

The meaning of the action requires clarification. Under the multiplicative interpretation suggested by the notation, for \(x\in F^\times=F\setminus\{0\}\) one defines
\[
x^{\sum_{g\in G}n_g g}=\prod_{g\in G}g(x)^{n_g}.
\]
This gives the following precise candidate for the intended hypothesis on nonzero elements:
\[
\forall x\in F^\times\quad\exists (n_g)_{g\in G}\in\Z^G\setminus\{(0)_{g\in G}\}
\quad \prod_{g\in G}g(x)^{n_g}\in S^\times.
\]
However, this multiplicative action is not defined at \(x=0\) for arbitrary group-ring elements with negative coefficients. The source explicitly quantifies over all of \(F\), so the displayed candidate cannot be asserted to be an exact transcription without resolving the convention at zero. There is also a natural additive group-ring action \(f\cdot x=\sum_{g\in G}n_g g(x)\), which yields a different hypothesis; the supplied record does not explicitly specify which action is intended.

Here \(F/S\) is purely inseparable if it is algebraic and, when \(\operatorname{char}(S)=p>0\), every \(a\in F\) satisfies \(a^{p^r}\in S\) for some integer \(r\geq0\) depending on \(a\). In characteristic zero, a purely inseparable extension is the trivial extension. The source further states that the analogous question for an uncountable field has an affirmative answer; this is a supplied background assertion, not independently verified here.
\subsection{Short English statement}
For a countably infinite field with finitely many specified automorphisms, suppose every element can be sent into a subfield using a nonzero integral group-ring element. Must the field equal that subfield or be purely inseparable over it? The action, especially its meaning at zero under the multiplicative interpretation, needs clarification before the supplied formulation is fully precise.
\subsection{Sources}
\begin{itemize}
\item Supplied local UnsolvedMath registry record, permanent ID \(700887\), exact source problem code \texttt{catalogue entry}, titled \emph{Kourovka Notebook Problem 12.69}. The record labels the problem \emph{open}; this is supplied registry status, not an independently verified literature assessment. Upstream revision: \texttt{d371ee917ac9647fbb06cd6fc529deec6f3dda82}.
\item Attributed to K. N. Ponomar\"ev, proposed in 1992, as Problem 12.69 of \emph{The Kourovka Notebook (Unsolved Problems in Group Theory)}, 12th issue (1992). The supplied background also identifies the 21st issue edition (2026), \texttt{arXiv:1401.0300v48}. The registry describes the collection as edited by E. I. Khukhro and V. D. Mazurov. These bibliographic details are supplied attribution and have not been independently checked.
\end{itemize}
\end{document}
